\section{Conclusion}
\label{sec:conclusion}

Communication verification asks whether interactions conform to a specification, state-machine verification whether transitions preserve invariants, and commitment verification whether obligations are managed. Dependency--Mechanism Adequacy asks whether the dependency a protocol exists for was managed, and whether a participant can establish that it was. The paper has shown that the second question factors into conduct, observability and interpretation, and that observability has an exact ceiling: no sound judge can claim more than its observations determine, and a canonical judge claims exactly that much. Epistemic adequacy is therefore a property of what a protocol makes observable.

Built as executable semantics over eight scenario families and thirty-six fault profiles, the calculus produced exact results. No single paradigm is safe across the families. A hybrid that refuses to accept a participant's report as evidence of its own performance, and refuses to read a record's silence beyond its watermark, made no false claim under any profile, and removing exactly those refusals made it unsafe in every family. False completion followed the evidence model, not the paradigm. Composition was decidable only with independent evidence of the remote party's stage. These conclusions held for arms written independently from the specifications, and the differential check that established it corrected the reference itself.

The result most useful to a designer is the ceiling. Before writing a judge, one can compute from a protocol's interface and an environment how much of a dependency's outcome any sound judge could ever establish, and so whether a richer interface, rather than a better judge, is what the dependency needs.
